\documentclass{article}
\usepackage{amsmath}
\usepackage{amsthm}
\usepackage{amsfonts}
\usepackage{amssymb}
\usepackage{bbold}
\usepackage{dsfont}
\usepackage{graphicx}
\usepackage{cleveref}
\usepackage{bm}
\usepackage{booktabs}
\usepackage{siunitx}
\usepackage{xcolor}
\usepackage{pdfpages}
\usepackage[utf8]{inputenc}
\usepackage[english]{babel}
\newtheorem{definition}{Definition}
\newtheorem{example}{Example}
\newcommand{\exP}[1]{\mathrm{e}^{#1}}
\newcommand{\dif}{\mathrm{d}}
\renewcommand{\vec}[1]{\mathbf{#1}}
\newcommand{\pdagger}{{\phantom{\dagger}}}
\newcommand{\state}[1]{\vert #1\rangle}
\newcommand{\cstate}[1]{\langle #1\vert}
\newcommand{\tr}[1]{\mathrm{tr}\left\{#1\right\}}
\newcommand{\trs}[2]{\, \mathrm{tr}_{#1}\{#2\}}
\newcommand{\expc}[3]{\langle #1\vert #2 \vert #3\rangle}
\usepackage[doi=false,url=false,isbn=false,backend=bibtex,style=nature,natbib=true,sorting=none]{biblatex}
\bibliography{bib.bib}
\usepackage{geometry}
\geometry{
	a4paper,
	total={170mm,257mm},
	left=15mm,
	right=15mm,
	top=15mm,
}
\begin{document}
\title{Ehrenfest-HEOM analysis of Fransson model}
\author{Rudolf Smorka}
\maketitle
\section{Model}
Fransson considers a model of a single-stranded helical polymer, within an atomistic approach. Each monomer is modeled as a single site. He considers, nearest-neighbor hopping, next nearest neighbor spin orbit coupling and vibrationally assisted hopping and SOC. The electronic part reads
\begin{equation}
	\begin{split}
		\hat{H}_\mathrm{el} &= \sum_{j,\sigma} \varepsilon_{j\sigma}^\pdagger \hat{d}_{j\sigma}^\dagger \hat{d}_{j\sigma}^\pdagger - t_0 \sum_{j\sigma} (\hat{d}_{j\sigma}^\dagger \hat{d}_{j+1,\sigma}^\pdagger + \mathrm{h.c.}) +\\& \lambda_0^\pdagger\sum_{j\sigma\sigma'} (i\hat{d}_{j\sigma}^\dagger (\vec{v}_j^{(+)}\cdot \bm{\sigma})_{\sigma\sigma'}^\pdagger \hat{d}_{j+2,\sigma'}^\pdagger + \mathrm{h.c.}),
	\end{split}
\end{equation}
and the coupling between electrons and vibrations
\begin{align}
	\hat{H}_\mathrm{el-ph} &= - t_1^\pdagger \sum_{j,\sigma}  \hat{d}_{j\sigma}^\dagger \hat{d}_{j+1\sigma}^\pdagger (\hat{a}_{j}^\dagger + \hat{a}_{j}^\pdagger) \\ & + \lambda_1^\pdagger\sum_{j,\sigma,\sigma'} (i \hat{d}_{j\sigma}^\dagger (\vec{v}_j^{(+)} \cdot \bm{\sigma})_{\sigma\sigma'}^\pdagger \hat{d}_{j+2\sigma'}^\pdagger +\mathrm{h.c.})(\hat{a}_{j}^\dagger + \hat{a}_{j}^\pdagger),\label{eq:HelPh}
\end{align}
and the vibrational part
\begin{equation}
	\hat{H}_\mathrm{ph} = \sum_{j} \omega_{j}^\pdagger \hat{a}_{j}^\dagger \hat{a}_{j}^\pdagger.
	\label{eq:HPh} 
\end{equation}

\section{Parameters considered in Franssons work}
\begin{align*}
	\varepsilon_0 &= -\SI{240}{\milli\electronvolt}\\
	\omega_0 &= \SI{50}{\milli\electronvolt}\\
	t_0 &= \SI{40}{\milli\electronvolt}\\
	t_1 &= 4,\SI{20}{\milli\electronvolt}\\
	\lambda_0 &= \SI{1}{\milli\electronvolt}\\
	\lambda_1 &= 0.1,\SI{1}{\milli\electronvolt}\\
	\Gamma &= \SI{10}{\milli\electronvolt}\\
	T &= \SI{25}{\milli\electronvolt}.
\end{align*}

\section{Analysis of spectrum of isolated molecule (exact diagonalization)}
\begin{itemize}
	\item Variation of static spin orbit coupling $\lambda_0$: Each eigenstate remains doubly degenerate. SOC part is invariant under $SU(2)$ trafos.
	\item Same holds true for variation of $\lambda_1$. 
	\item Variation of vibrational coordinates $X$ (classical), leads to unequal energy differences of eigenstates. 
	\item One way of looking at the Hamiltonian within the Ehrenfest approach, is that it acquires a time dependence through the time-evolution of $\{X_i(t)\}$. 
\end{itemize}
\section{Analysis of spectrum of coupled molecule and time-independent transport (without vibrations, using NEGF)}
\begin{itemize}
	\item Molecule is coupled to the $P$-polarized left lead via $\Gamma_L = \Gamma_0 (\mathds{1}+P\sigma_z)$ and to the unpolarized right lead via $\Gamma_R=\Gamma_0$. 
	\item Lead polarization $P$ reduces/enhances transmission function of certain states. These states are typically with the "correct" spin localized near the left interface. 
	\item Although transmission function shows differences between $P=+1$ and $P=-1$, they are tiny for the parameters considered by Fransson and their effect on the current-voltage characteristics is basically negligible.
	\item In summary: Spin selectivity in the model without vibrations is negligibly small. 
\end{itemize}
\section{Analysis of dynamics and spin selectivity in the full model within Ehrenfest-HEOM}
\begin{itemize}
	\item For the parameters above, stationary state is reached after $>\SI{50}{\pico\second}$. Current fluctuates in time. Could be finite sampling artifact, but I observe that the current fluctuates even in absence of interactions for a very long time. Spin selectivity, since the charge currents are time-dependent, is also time-dependent. At some times, it is negative, and even changes sign. For the system sizes $L=2\times 4$ and  $L=3\times 4$, I have data that looks somewhat converged. There, spin selectivity is no greater than $S=0.5\%$.  
	\item Variation of $\lambda_0$: I have data for $\lambda_0 = \SI{1}{\milli\electronvolt}$ and $\lambda_0 = \SI{10}{\milli\electronvolt}$. Does not help increasing spin selectivity.
	\item Variation of $\lambda_1$: I have data for $\lambda_1 = \SI{0.1}{\milli\electronvolt}$, $\lambda_1 = \SI{5}{\milli\electronvolt}$ and $\lambda_1 = \SI{10}{\milli\electronvolt}$. For $\lambda_1 = \SI{5}{\milli\electronvolt}$, spin selectivity is somewhat increased, reaching up to $S\sim 1\%$. For $\SI{10}{\milli\electronvolt}$, the picture is not so clear because I was not able to reach a steady state yet. Calculations for longer times are still running.
	\item Voltage: Looks like for large voltage (when basically all eigenstates are in the bias window) spin selectivity is reduced to zero. It looks like it is maximal when the voltage is such that less or equal to half of the states are in the bias window.
	\item Thought on parameters: $\varepsilon_0=-\SI{240}{\milli\electronvolt}$ seems to play a big role as well. It makes the results hard to converge, and for this parameter spin selectivity seems to be particularly large. I also tried $\varepsilon_0 = \SI{0}{\electronvolt}$, and the spin selectivity dropped to 0. 
	\item I don't understand the effect of the vibrations yet. Also, I don't understand how the vibrations affect the electrons, such that there is a spin selectivity in the end. 
	\item Overall, it seems like the spin-asymmetric $\Gamma_L$, SOC, vibrations, and bound states play a cooperative role in creating a spin selectivity.
\end{itemize}

\end{document}


